\documentclass[11pt]{article}

\usepackage[T1]{fontenc}
\usepackage{amsmath,amssymb,amsthm}
\usepackage{hyperref}

\title{Ordinal Profile-LRT Finite-Sample Calibration}
\author{}
\date{\today}

\newcommand{\E}{\mathbb{E}}
\newcommand{\Var}{\operatorname{Var}}
\newcommand{\Cov}{\operatorname{Cov}}
\newcommand{\argmin}{\operatorname*{arg\,min}}
\newcommand{\plim}{\operatorname*{plim}}

\theoremstyle{plain}
\newtheorem{definition}{Definition}
\newtheorem{proposition}{Proposition}
\theoremstyle{remark}
\newtheorem*{remark}{Remark}

\begin{document}
\maketitle

\section*{Purpose}

This note narrows the profile-LRT small-sample problem to all-ordinal
limited-information SEM. The target is not a universal Bartlett correction.
The target is a decision rule for profile-LRT confidence intervals for
polychoric omega under DWLS first, with WLS as a later comparison, where
thresholds, polychorics, NACOV, and fitted weights are estimated from
categorical margins.

The companion note \cite{profilebartlettnote} records the general distinction:
when the limiting reference is a weighted chi-square spectrum, a scalar
Bartlett factor is only a first-moment correction. Here the focal profile
restriction has one degree of freedom. After robust or misspecification scaling,
the leading reference again has the form $\chi^2_1$, so a finite-sample scalar is
at least meaningful. The question is whether it is useful, stable, and
diagnostically separable from ordinal point-estimate bias.

\section{Object}

Let $s$ collect ordinal sample thresholds and polychoric correlations, and let
$F(\theta;s,W)$ be the limited-information DWLS or WLS criterion. For a scalar
functional $g(\theta)$, here polychoric/latent-response omega, define
\[
T(g_0) =
  N\{F(\tilde\theta(g_0);s,W) - F(\hat\theta;s,W)\},
\qquad
\tilde\theta(g_0)=\argmin_{\theta:g(\theta)=g_0} F(\theta;s,W).
\]
The current magmaan profile helpers also compute one-degree-of-freedom reference
scalings at the constrained endpoint. Write the reference-scaled statistic as
\[
Z(g_0) =
\begin{cases}
T(g_0), & \text{ordinary},\\
T(g_0)/\lambda_{\mathrm{rob}}, & \text{robust scaled},\\
T(g_0)/\lambda_{\mathrm{mis}}, & \text{misspecification scaled}.
\end{cases}
\]
The asymptotic reference for $Z(g_0)$ is $\chi^2_1$. The finite-sample question
is whether
\[
  Z(g_0) \le c\,\chi^2_{1,1-\alpha}
\]
with a stable $c$ improves coverage in the small ordinal cells.

\section{Bias Decomposition}

The ordinal problem has four separable failure modes.

\begin{enumerate}
\item \emph{LR inflation.} Even after the robust endpoint scaling, $\E Z$ may
exceed one at small $N$. This is the only component a Bartlett-type threshold
factor directly targets.
\item \emph{Reference-scale error.} The endpoint $\lambda$ is itself estimated
from thresholds, polychorics, NACOV, and sometimes fitted weights. Sparse cells
can move the robust or misspecification reference before any Bartlett factor is
considered.
\item \emph{Functional point bias.} The unrestricted omega estimate may be
biased because the polychoric stage and the fitted CFA surface are biased at
small $N$. A threshold correction can cover the population value at $g_0$, but
it cannot make a biased point estimate a good center for width summaries.
\item \emph{Boundary and feasibility mass.} For bounded reliability coefficients,
some samples may put the profile surface near a boundary or make $g_0$ hard to
attain. In the maximal-reliability experiments this mass dominated the mean
inflation and defeated analytic Lawley and per-sample bootstrap factors.
\end{enumerate}

\section{Candidate Corrections}

\paragraph{Oracle cell constant.}
For a design cell, estimate $c_{\mathrm{oracle}}=\E Z(g_{\mathrm{pop}})$ by
Monte Carlo and report coverage under
$Z(g_{\mathrm{pop}})\le c_{\mathrm{oracle}}\chi^2_{1,1-\alpha}$. This is not a
method, but it identifies whether the statistic is repairable by a scalar.

\paragraph{Ordinal null-bootstrap constant.}
For a design class, simulate ordinal data from the constrained fitted model
imposing $g=g_0$, refit the unrestricted and constrained models, and average the
reference-scaled $Z$. The value is used as a cell-level constant, not as a
per-sample random factor. The earlier maximal-reliability result rules out
per-sample use: the bootstrap factor was anti-correlated with the samples that
needed correction most.

\paragraph{Shrinkage constant.}
If the null-bootstrap constant is noisy but directionally useful, shrink it to a
grid prior:
\[
  c_{\mathrm{shrink}} = w c_{\mathrm{boot}} + (1-w)c_{\mathrm{grid}}.
\]
The weight $w$ should be chosen by coverage and anti-correlation diagnostics, not
by mean recovery alone.

\paragraph{Weight and NACOV regularization.}
If reference-scale error dominates, the correction target is not $c$ but the
estimated ordinal $\Gamma$ and $W$. DLS, EBA/pEBA-style spectrum shrinkage, or
regularized NACOV can be tested as inputs to the profile reference. These are
finite-sample reference corrections, not Bartlett factors.

\paragraph{Bias-aware reporting.}
If point-functional bias dominates, report the profile interval as calibrated for
test inversion but keep a separate bias diagnostic for omega. Do not hide
functional bias inside the Bartlett factor.

\section{Experiment 53}

Experiment 53 is the first taxonomy run. It uses all-ordinal one-factor probit
data, DWLS fits, and the existing
\texttt{frontier\_profile\_lrt\_ordinal\_polychoric\_omega} surface. It records
ordinary, robust-scaled, and misspecification-scaled profile tests at the
population omega, plus the fit-free robust delta interval for comparison. The
initial grid deliberately defers the WLS robust-scaled arm to spend the
simulation budget on DWLS replications.

The decision rule is:
\begin{enumerate}
\item If uncorrected $Z$ is near $\chi^2_1$ and delta coverage is acceptable,
there is no ordinal Bartlett problem in that cell.
\item If $Z$ is inflated and the oracle constant fixes coverage, advance to an
ordinal null-bootstrap constant.
\item If $Z$ is not inflated but coverage is poor, prioritize omega point-bias
and sparse-threshold diagnostics.
\item If robust and misspecification references diverge materially, keep the
small-sample constant conditional on the selected reference regime.
\end{enumerate}

The first full DWLS-only grid used 1000 replications per cell over
$N\in\{50,100,250,500\}$ and balanced versus threshold-extreme cuts. Ordinary
DWLS profile inversion is not a viable reference here: coverage stayed around
0.60--0.69 because the unscaled statistic has mean roughly 3.5--6.4. Endpoint
robust scaling and the misspecification-scaled reference both move the reference
coverage close to nominal in most cells, roughly 0.95--0.97. The exception is
the threshold-extreme $N=50$ cell, where empty or near-empty categories create
profile failures and unstable fitted-weight reference scales; coverage is about
0.93 rather than 0.95. The practical next correction target is therefore sparse
ordinal-summary / fitted-weight stability in the small extreme-threshold cell,
not a broad Bartlett correction for every ordinal omega interval.

The misspecified extension adds a locally dependent latent-response DGP, with
residual correlations between items 1--2 and 5--6, while fitting the same
one-factor DWLS model. The target is the large-sample pseudo-true
model-implied profile omega from the working one-factor fit, not the generating
model's reliability. With 1000 replications per cell and pseudo-targets from
100000 generated cases, ordinary DWLS again undercovers badly
(about 0.59--0.69). Endpoint robust scaling covers about 0.94--0.97, and the
misspecification-scaled reference covers about 0.94--0.97. The weakest
misspecified cell remains threshold-extreme $N=50$, with many profile failures
and unstable fitted-weight scales. Thus the misspecification reference is
behaving as intended for the working-model target, but it does not remove the
sparse ordinal-summary problem.

The scalar misspecification reference should be read as scaled, not as a
substantive mixture distinction. For a one-degree-of-freedom profile test the
weighted-\(\chi^2\) law has a single eigenvalue \(\lambda\), so comparing
\(T\) to \(\lambda\chi^2_1\) is equivalent to comparing \(T/\lambda\) to
\(\chi^2_1\). The experiment keeps any older ``mixture'' spelling only as a
legacy alias for this scalar reference.

The stress extension adds two harder working-model-misspecification cells:
stronger residual local dependence and a two-factor population, both still fit
by the same one-factor DWLS model. The stress default uses
\(N\in\{25,35,50,75\}\) and threshold-extreme / threshold-sparse cuts. The
sparse-threshold cells mostly diagnose ordinal summary construction rather than
reference-law calibration: empty categories dominate the failures. Among cells
with at least 200 successful profile fits in the 300-replication pointwise run,
the useful hard cases are threshold-extreme \(N=50\). Robust-scaled coverage was
0.932 under strong local dependence and 0.926 under the two-factor DGP; the
oracle-constant diagnostic moved these to 0.964 and 0.961, respectively. The
misspecification-scaled arm showed the same scalar-calibration signal in the
strong local-dependence cell (0.935 to 0.960) and a smaller one in the two-factor
cell (0.948 to 0.967). These are the first cells where a finite-sample scalar
constant is plausibly worth testing.

The first CI-inversion pass used the existing profile-CI root finder with
200 replications per cell and the same pseudo-targets. Robust-scaled inversion
is the stable interval lane: among 3168 attempted robust-scaled CIs, only 38
failed, all target-test/CI inclusion decisions agreed, and cell coverage ranged
from about 0.918 to 0.975. Misspecification-scaled inversion is not yet a
stable CI method: among 3154 attempted CIs under the legacy singleton-mixture
spelling, 1201 failed, mostly because the endpoint robust profile quadratic was
not positive. Successful misspecification-scaled CIs cover near nominal, but
that is a selected subset. Thus the current connected-CI evidence supports
robust-scaled inversion; the misspecification-scaled path is still a pointwise
reference diagnostic until endpoint reference construction is hardened.

A targeted stress CI pass on threshold-extreme \(N=50\), with 100 replications
per hard DGP, preserves the same split. Robust-scaled CIs remain connected: 7 of
92 strong-local attempts and 18 of 94 two-factor attempts failed, successful
coverage was 0.929 and 0.934, and the pointwise-test / CI inclusion decisions
never disagreed. Misspecification-scaled endpoint inversion is still not a
usable interval lane in these sparse hard cells: 178 of 179 attempted CIs failed,
primarily because the endpoint robust profile sandwich quadratic was not
positive. The immediate CI target is therefore robust-scaled inversion plus a
candidate scalar finite-sample constant in the \(N=50\) threshold-extreme stress
cells.

\section*{Working Verdict}

The plausible contribution is not ``a Bartlett correction for categorical SEM''.
The plausible contribution is a finite-sample calibration taxonomy for ordinal
profile-LRT reliability intervals. A scalar constant remains on the table only
when the reference-scaled one-degree-of-freedom statistic is the source of
miscoverage. Otherwise the correction belongs to the ordinal summary stage or to
the reliability target itself.

\begin{thebibliography}{99}

\bibitem{profilebartlettnote}
Magmaan note. Profile likelihood-ratio intervals, the Bartlett correction, and
the eigenvalue-spectrum reference law.
\texttt{docs/research/notes/profile\_lrt\_bartlett\_spectral.tex}.

\bibitem{wilks1938}
Wilks, S. S. (1938). The large-sample distribution of the likelihood ratio for
testing composite hypotheses. \emph{Annals of Mathematical Statistics}, 9(1),
60--62.
\href{https://doi.org/10.1214/aoms/1177732360}{doi:10.1214/aoms/1177732360}.

\bibitem{bartlett1937}
Bartlett, M. S. (1937). Properties of sufficiency and statistical tests.
\emph{Proceedings of the Royal Society of London A}, 160(901), 268--282.
\href{https://doi.org/10.1098/rspa.1937.0109}{doi:10.1098/rspa.1937.0109}.

\bibitem{satorra1989}
Satorra, A. (1989). Alternative test criteria in covariance structure analysis:
a unified approach. \emph{Psychometrika}, 54(1), 131--151.
\href{https://doi.org/10.1007/BF02294453}{doi:10.1007/BF02294453}.

\bibitem{hallinoue2003}
Hall, A. R., and Inoue, A. (2003). The large sample behaviour of the generalized
method of moments estimator in misspecified models. \emph{Journal of
Econometrics}, 114(2), 361--394.
\href{https://doi.org/10.1016/S0304-4076(03)00089-7}{doi:10.1016/S0304-4076(03)00089-7}.

\bibitem{lee2014}
Lee, S. (2014). Asymptotic refinements of a misspecification-robust bootstrap
for generalized method of moments estimators. \emph{Journal of Econometrics},
178, 398--413.
\href{https://doi.org/10.1016/j.jeconom.2013.05.008}{doi:10.1016/j.jeconom.2013.05.008}.

\bibitem{yuantianyanagihara2015}
Yuan, K.-H., Tian, Y., and Yanagihara, H. (2015). Empirical correction to the
likelihood ratio statistic for structural equation modeling with many variables.
\emph{Psychometrika}, 80(2), 379--405.
\href{https://doi.org/10.1007/s11336-013-9386-5}{doi:10.1007/s11336-013-9386-5}.

\end{thebibliography}

\end{document}
